\documentclass[11pt,a4paper]{article}

\usepackage[T1]{fontenc}
\usepackage[utf8]{inputenc}
\usepackage{lmodern}
\usepackage{amssymb,amsmath,amsthm,mathtools}
\usepackage[margin=2.5cm]{geometry}
\usepackage{parskip}
\usepackage[hidelinks]{hyperref}
\usepackage{microtype}

% Theorem environments (arXiv / journal style)
\theoremstyle{plain}
\newtheorem{theorem}{Theorem}[section]
\newtheorem{proposition}[theorem]{Proposition}
\newtheorem{lemma}[theorem]{Lemma}
\newtheorem{corollary}[theorem]{Corollary}

\theoremstyle{definition}
\newtheorem{definition}[theorem]{Definition}
\newtheorem{example}[theorem]{Example}
\newtheorem{remark}[theorem]{Remark}

% Notation
\newcommand{\Type}{\mathbf{Type}}
\newcommand{\id}{\mathrm{id}}
\newcommand{\Cont}{\mathbf{Cont}}
\newcommand{\Shape}{\mathrm{Shape}}
\newcommand{\Pos}{\mathrm{Pos}}
\newcommand{\seq}{\mathbin{\triangleleft}}

\title{The Sequential Monoidal Structure on Containers}
\author{Extracted from a machine-checked development}
\date{}

\begin{document}
\maketitle

\begin{abstract}
The category $\Cont$ of containers and container morphisms carries a
monoidal structure whose tensor is the sequential (composition)
operator~$\seq$ and whose unit is the identity container~$I$. We record
the bifunctoriality of~$\seq$ on morphisms, the naturality of the
associator and unitors, and the two coherence laws---the pentagon and the
triangle. All equalities are transport-free: position fibres line up
definitionally, so every law holds by direct computation. The proofs of
all stated results have been machine-checked; we record only the
mathematical content.
\end{abstract}

\tableofcontents
\bigskip

%======================================================================
\section{Introduction}
%======================================================================

Containers, in the sense of Abbott, Altenkirch, and Ghani, form a
category $\Cont$ whose objects are pairs $(S,P)$ of a shape type and a
position family, and whose morphisms are shape maps with contravariant
position maps. Composition of the associated polynomial functors lifts to
a tensor product on containers: the \emph{sequential operator}~$G \seq F$
puts an $F$-structure at every position of a $G$-structure. Its unit is
the identity container~$I$ (one shape, one position).

The structural isomorphisms---associator and left/right unitors---are
standard and transport-free. What remains to certify $(\Cont, \seq, I)$
as a monoidal category is bifunctoriality of~$\seq$ on morphisms and the
two Mac Lane coherence laws. That is the content of this note.

\paragraph{What this note claims.}
In ordinary mathematical language: $\seq$ is a bifunctor; the associator
and unitors are natural; the pentagon and triangle commute. On shapes the
associator is currying; on positions it is $\Sigma$-reassociation. Both
legs of the pentagon flatten the fourfold composite
$((C_1 \seq C_2) \seq C_3) \seq C_4$ to the right-nested normal form
$C_1 \seq (C_2 \seq (C_3 \seq C_4))$, and on positions every map is the
identity on the underlying flat tuple.

We work throughout in a single universe of types. Proofs are indicated by
the phrase ``by lean.''

%======================================================================
\section{Background}
%======================================================================

\begin{definition}[Container]\label{def:container}
A \textbf{container} $C$ consists of a type $\Shape(C)$ of shapes and,
for each shape~$s$, a type $\Pos(C)(s)$ of positions.
\end{definition}

\begin{definition}[Container morphism]\label{def:morphism}
A \textbf{morphism} $\mu : C \to D$ consists of a shape map
$u : \Shape(C) \to \Shape(D)$ and, for each $s$, a position map
$\Pos(D)(u(s)) \to \Pos(C)(s)$. Composition and identities are the
evident ones; we write $\Cont$ for the resulting category.
\end{definition}

\begin{definition}[Sequential operator]\label{def:seq}
The \textbf{sequential composite} $G \seq F$ of containers $G$ and $F$
has
\begin{align}
  \Shape(G \seq F)
  &\;=\;
  \sum_{t \,:\, \Shape(G)}
  \bigl(\Pos(G)(t) \to \Shape(F)\bigr), \\[4pt]
  \Pos(G \seq F)\bigl\langle t,\, f \bigr\rangle
  &\;=\;
  \sum_{p \,:\, \Pos(G)(t)}
  \Pos(F)\bigl(f(p)\bigr).
\end{align}
An element of the extension $\mathrm{Ext}\,(G \seq F)\,X$ is a
$G$-structure of $F$-structures of $X$-labels---equivalently,
$\mathrm{Ext}\,G\bigl(\mathrm{Ext}\,F\,X\bigr)$.
\end{definition}

\begin{definition}[Unit container]\label{def:I}
The \textbf{identity container} $I$ has a single shape and a single
position at that shape. It is the unit for~$\seq$.
\end{definition}

\begin{remark}[Structural isomorphisms]
The associator
\[
  \alpha_{X,Y,Z}
  \;:\;
  (X \seq Y) \seq Z
  \;\xrightarrow{\;\sim\;}\;
  X \seq (Y \seq Z)
\]
acts by currying on shapes and by $\Sigma$-reassociation on positions.
The left unitor $\lambda_X : I \seq X \xrightarrow{\sim} X$ and right
unitor $\rho_X : X \seq I \xrightarrow{\sim} X$ are the evident
collapses. All three are isomorphisms of containers (transport-free),
and we take them as given from the sequential development.
\end{remark}

%======================================================================
\section{Bifunctoriality}
%======================================================================

\begin{definition}[Action on morphisms]\label{def:seq2}
Given morphisms $\mu : G \to G'$ and $\nu : F \to F'$, the
\textbf{sequential action}
\[
  \mu \seq_2 \nu
  \;:\;
  G \seq F \;\longrightarrow\; G' \seq F'
\]
is the container morphism with
\begin{align}
  \text{on shapes:}\quad
  \bigl\langle t,\, f \bigr\rangle
  &\;\longmapsto\;
  \bigl\langle
    \mu(t),\;
    \nu \circ f \circ \mu^\sharp_t
  \bigr\rangle, \\[4pt]
  \text{on positions:}\quad
  \bigl\langle p,\, q \bigr\rangle
  &\;\longmapsto\;
  \bigl\langle
    \mu^\sharp_t(p),\;
    \nu^\sharp_{f(\mu^\sharp_t(p))}(q)
  \bigr\rangle,
\end{align}
where $\mu^\sharp$ and $\nu^\sharp$ denote the (contravariant) position
components of $\mu$ and $\nu$. The position fibres of source and target
agree definitionally, so no transport appears.
\end{definition}

\begin{definition}[Whiskering]\label{def:whisker}
\begin{enumerate}
  \item \textbf{Left whiskering}
        $G \seq \nu : G \seq F \to G \seq F'$ is
        $\id_G \seq_2 \nu$.
  \item \textbf{Right whiskering}
        $\mu \seq F : G \seq F \to G' \seq F$ is
        $\mu \seq_2 \id_F$.
\end{enumerate}
\end{definition}

\begin{proposition}[Functor laws]\label{prop:seq2-laws}
For all containers and morphisms of the appropriate types:
\begin{enumerate}
  \item $\id_G \seq_2 \id_F = \id_{G \seq F}$;
  \item $(\mu \mathbin{;}\mu') \seq_2 (\nu \mathbin{;}\nu')
        = (\mu \seq_2 \nu) \mathbin{;} (\mu' \seq_2 \nu')$
        (interchange).
\end{enumerate}
\end{proposition}

\begin{proof}
By lean.
\end{proof}

Thus $\seq$ is a bifunctor $\Cont \times \Cont \to \Cont$.

%======================================================================
\section{Naturality of the structural isomorphisms}
%======================================================================

\begin{proposition}[Naturality of the associator]\label{prop:assoc-nat}
For morphisms $f_i : X_i \to Y_i$ ($i=1,2,3$),
\begin{equation}
  \bigl((f_1 \seq_2 f_2) \seq_2 f_3\bigr)
  \mathbin{;}
  \alpha_{Y_1,Y_2,Y_3}
  \;=\;
  \alpha_{X_1,X_2,X_3}
  \mathbin{;}
  \bigl(f_1 \seq_2 (f_2 \seq_2 f_3)\bigr).
\end{equation}
\end{proposition}

\begin{proof}
By lean.
\end{proof}

\begin{proposition}[Naturality of the unitors]\label{prop:unit-nat}
For every morphism $f : X \to Y$,
\begin{align}
  (I \seq f) \mathbin{;} \lambda_Y
  &\;=\;
  \lambda_X \mathbin{;} f, \\[4pt]
  (f \seq I) \mathbin{;} \rho_Y
  &\;=\;
  \rho_X \mathbin{;} f.
\end{align}
\end{proposition}

\begin{proof}
By lean.
\end{proof}

\begin{remark}
Each square holds because both legs are the same combination of
currying/$\Sigma$-reassociation with the morphism data threaded
identically.
\end{remark}

%======================================================================
\section{Coherence: pentagon and triangle}
%======================================================================

\begin{theorem}[Pentagon identity]\label{thm:pentagon}
For all containers $C_1,C_2,C_3,C_4$, the two ways of reassociating
$((C_1 \seq C_2) \seq C_3) \seq C_4$ to
$C_1 \seq (C_2 \seq (C_3 \seq C_4))$ coincide:
\begin{equation}
\begin{aligned}
  &\bigl(
    (\alpha_{C_1,C_2,C_3} \seq C_4)
    \mathbin{;}
    \alpha_{C_1,\, C_2 \seq C_3,\, C_4}
  \bigr)
  \mathbin{;}
  (C_1 \seq \alpha_{C_2,C_3,C_4})
  \\
  &\;=\;
  \alpha_{C_1 \seq C_2,\, C_3,\, C_4}
  \mathbin{;}
  \alpha_{C_1,\, C_2,\, C_3 \seq C_4}.
\end{aligned}
\end{equation}
\end{theorem}

\begin{proof}
By lean.
\end{proof}

\begin{remark}
Both legs flatten the left-nested fourfold composite to the right-nested
normal form. On positions, every map is the identity on the underlying
flat tuple of positions, so the two paths agree on the nose.
\end{remark}

\begin{theorem}[Triangle identity]\label{thm:triangle}
For all containers $C_1,C_2$,
\begin{equation}
  \alpha_{C_1,\, I,\, C_2}
  \mathbin{;}
  (C_1 \seq \lambda_{C_2})
  \;=\;
  \rho_{C_1} \seq C_2,
\end{equation}
as morphisms $(C_1 \seq I) \seq C_2 \to C_1 \seq C_2$.
\end{theorem}

\begin{proof}
By lean.
\end{proof}

%======================================================================
\section{The monoidal category}
%======================================================================

\begin{definition}[Isomorphism]\label{def:iso}
An \textbf{isomorphism} $X \cong Y$ in a category is a morphism with a
two-sided inverse.
\end{definition}

\begin{definition}[Monoidal category]\label{def:monoidal}
A \textbf{monoidal category} consists of a category $\mathcal{C}$
together with
\begin{enumerate}
  \item a tensor product $\otimes : \mathcal{C} \times \mathcal{C} \to
        \mathcal{C}$ (an object map and a morphism map satisfying the
        two functor laws),
  \item a unit object $I$,
  \item natural isomorphisms
        $\alpha_{X,Y,Z} : (X \otimes Y) \otimes Z
        \xrightarrow{\sim} X \otimes (Y \otimes Z)$,
        $\lambda_X : I \otimes X \xrightarrow{\sim} X$,
        $\rho_X : X \otimes I \xrightarrow{\sim} X$,
  \item the pentagon and triangle coherence laws.
\end{enumerate}
\end{definition}

\begin{theorem}[$(\Cont, \seq, I)$ is monoidal]\label{thm:cont-monoidal}
The category of containers, equipped with the sequential operator as
tensor, the identity container as unit, and the structural isomorphisms
recalled above, is a monoidal category. Bifunctoriality is
Proposition~\ref{prop:seq2-laws}; naturality is
Propositions~\ref{prop:assoc-nat} and~\ref{prop:unit-nat}; coherence is
Theorems~\ref{thm:pentagon} and~\ref{thm:triangle}.
\end{theorem}

\begin{proof}
By lean.
\end{proof}

\begin{remark}
Every law holds by direct computation: the development is transport-free
and axiom-clean. In particular there is no appeal to Mac Lane's coherence
theorem as a black box; the pentagon and triangle are checked on the
nose for this concrete tensor.
\end{remark}

%======================================================================
\section{Discussion}
%======================================================================

\paragraph{What has been established.}
The sequential operator makes $\Cont$ into a monoidal category. The
structure is the spine of the broader development: comonoids for
$(\Cont, \seq, I)$ are precisely directed containers, so the monoidal
data certified here underwrite that correspondence.

\paragraph{Relation to other monoidal structures.}
The category $\Cont$ carries at least four monoidal structures of
interest (coproduct, product, Dirichlet tensor, and sequential
composition). Only one can be the canonical monoidal instance without
colliding in a formal system; the sequential structure is preferred here
because it is the composition of polynomial functors. The other three are
treated in companion notes.

\paragraph{Machine checking.}
Every definition and theorem statement above corresponds to a declaration
in a machine-checked formalisation. The phrase ``by lean'' indicates that
the proof has been verified by a proof assistant.

%======================================================================
\section*{Acknowledgements}
%======================================================================

The mathematical content follows the theory of containers of Abbott,
Altenkirch, and Ghani, and the monoidal structure of sequential
composition of polynomial functors. The presentation here is a prose
rendering of a formal development intended for human review of the
theorem statements.

\begin{thebibliography}{9}

\bibitem{abbott2005}
M.~Abbott, T.~Altenkirch, and N.~Ghani.
\newblock Containers: constructing strictly positive types.
\newblock \emph{Theoretical Computer Science}, 342(1):3--27, 2005.

\bibitem{abbott2003}
M.~Abbott, T.~Altenkirch, and N.~Ghani.
\newblock Categories of containers.
\newblock In \emph{FoSSaCS}, volume~2620 of \emph{LNCS}, pages 23--38.
\newblock Springer, 2003.

\bibitem{maclane1971}
S.~Mac~Lane.
\newblock \emph{Categories for the Working Mathematician}.
\newblock Springer, 1971.

\end{thebibliography}

\end{document}
